@@M0@@
By Proposition 6.11, the homomorphism: @@M1@@ is an isomorphism if and only if the homomorphism: @@M2@@ is one. Since it is visibly surjective, the latter is also surjective, and it therefore suffices to show that @@M3@@ is injective. By functoriality of @@M4@@, we obtain the diagram:
@@M5@@
since @@M6@@ has bounded degrees, the right vertical arrow is an isomorphism (6.12), and using 6.11, we therefore obtain a morphism: @@M7@@, the composite of the two being the identity by the universal property of @@M8@@.

\sgaStableTarget{sga6.e05.sec.7}\section*{APPENDIX: The Adams operations @@M9@@}

We follow here a presentation due to J.P. Serre.

\sgaStableTarget{sga6.e05.para.7-1}7.1. Let @@M10@@ be a pre-@@M11@@-ring. For every @@M12@@, we denote by @@M13@@ the derivative with respect to @@M14@@ of @@M15@@. We then define elements of @@M16@@, denoted by @@M17@@, by putting:
\sgaStableTarget{sga6.e05.eq.7-1-1}@@M18@@

The maps @@M19@@ are additive; indeed, we have:
@@M20@@
